\section{Manfaat Penelitian}

Manfaat yang diharapkan dari penelitian ini adalah:
\begin{enumerate}
    \item \textbf{Kontribusi Teoretis}: Memberikan kontribusi dalam kajian relativitas umum dan fisika lubang hitam, khususnya dalam memahami perilaku perturbasi medan skalar pada dua latar ruang-waktu yang berbeda.
    \item \textbf{Edukasi}: Menyediakan referensi akademik yang sistematis dan analitis di lingkungan FSAD-ITS mengenai penurunan dan analisis perturbasi medan skalar dalam dua geometri ruang-waktu.
    \item \textbf{Landasan Riset Lanjutan}: Menjadi pijakan awal bagi studi perturbasi yang lebih kompleks, baik dengan mempertimbangkan jenis medan lain (elektromagnetik atau gravitasi), maupun jenis singularitas ruang-waktu yang lain seperti Kerr, Reissner-Nordstr\"{o}m, dan Kerr-Newman.
\end{enumerate}